\documentclass[11pt]{article}
\usepackage{amsmath,amsthm,amssymb}
\usepackage[margin=1in]{geometry}

\newtheorem{theorem}{Theorem}
\newtheorem{lemma}[theorem]{Lemma}
\newtheorem{proposition}[theorem]{Proposition}
\newtheorem{corollary}[theorem]{Corollary}
\newtheorem{conjecture}[theorem]{Conjecture}
\theoremstyle{definition}
\newtheorem{definition}[theorem]{Definition}
\theoremstyle{remark}
\newtheorem{remark}[theorem]{Remark}

\newcommand{\ZZ}{\mathbb{Z}}
\newcommand{\CC}{\mathbb{C}}
\newcommand{\Sym}{\mathfrak{S}}
\newcommand{\im}{\operatorname{im}}
\newcommand{\tr}{\operatorname{tr}}
\newcommand{\rk}{\operatorname{rank}}

\title{The Rank of the Symmetrizing Product: \\
Injectivity, Cascade Surjectivity, and the Formula $n!/2^{\lfloor n/2 \rfloor}$}
\author{Clio}
\date{April 14, 2026}

\begin{document}
\maketitle

\begin{abstract}
We study the staircase symmetrizing product
$\Pi = \prod(1+s_{i_j})/2$ for the longest element $w_0 \in \Sym_n$,
acting on the left regular representation. We prove an
\emph{Injectivity Lemma} showing rank drops occur only at the $P_1$
steps, a \emph{First Block Surjectivity} theorem identifying the image
after the first pass as $+1(s_{n-1})$, a \emph{Cascade Surjectivity}
lemma showing how eigenspace conditions propagate through the cascade,
and a \emph{Trace Preservation} lemma controlling $s_1$-traces. These
combine to prove the first two halvings unconditionally and yield the
rank formula $\rk(\Pi) = n!/2^{\lfloor n/2\rfloor}$ modulo a halving
lemma verified for $n \le 8$.
\end{abstract}


\section{Setup}

Let $\Sym_n$ act on $\{1, \dots, n\}$ with simple transpositions
$s_i = (i, i{+}1)$. In the left regular representation, $s_i$ acts by
$s_i \cdot e_w = e_{s_i w}$. Define projections
$P_i = (1 + s_i)/2$, which are idempotents projecting onto the
$+1$ eigenspace of~$s_i$.

\begin{definition}[Staircase product]
The staircase reduced word for $w_0 \in \Sym_n$ consists of $n{-}1$
\emph{passes}: pass~$j$ (for $j = 1, \dots, n{-}1$) applies the projections
$P_1, P_2, \dots, P_{n-j}$ in sequence. The full product is
\[
  \Pi = \underbrace{P_1}_{\text{pass }n{-}1}
  \cdot \underbrace{P_2 P_1}_{\text{pass }n{-}2}
  \cdots
  \underbrace{P_{n-1} \cdots P_2 P_1}_{\text{pass }1}.
\]
The rightmost pass is applied first.
\end{definition}

\noindent\textbf{Notation.} Write $+1(s_k)$ for the $+1$ eigenspace
of~$s_k$ and $-1(s_k)$ for the $-1$ eigenspace. Write $W_j$ for the
image of $\CC[\Sym_n]$ after the first $j$ passes have been applied.


\section{The Injectivity Lemma}

\begin{lemma}[Triple Lemma]
\label{lem:triple}
In $\CC[\Sym_n]$ with $n \ge 3$, for any $1 \le k \le n-2$:
\[
  \ker(P_{k+1}) \cap \im(P_k) = \{0\}.
\]
Equivalently, $+1(s_k) \cap -1(s_{k+1}) = \{0\}$.
\end{lemma}

\begin{proof}
By symmetry, it suffices to prove the case $k = 1$. Write
$v = \sum c_w \, e_w$ and group terms by the triple of values
$\{a, b, c\}$ occupying positions $1, 2, 3$ and by the tail
$w|_{\{4,\dots,n\}}$. Within each 6-element group, the condition
$s_1 v = v$ forces equal coefficients on each $s_1$-orbit, leaving
three variables $\alpha, \beta, \gamma$. The condition $s_2 v = -v$
gives $\beta = -\alpha$, $\gamma = -\alpha$, and $\gamma = -\beta$.
Together: $\gamma = -\gamma$, so $\alpha = \beta = \gamma = 0$.
\end{proof}

\begin{theorem}[Within-Pass Injectivity]
\label{thm:injectivity}
In the sequential application $P_1, P_2, \dots, P_m$ (with $m \le n-1$),
$P_j$ is injective on $P_{j-1}(\cdots P_1(W)\cdots)$ for all $j \ge 2$
and any initial subspace~$W$.
\end{theorem}

\begin{proof}
At step $j \ge 2$, the image of $P_{j-1}$ lies in $\im(P_{j-1}) = +1(s_{j-1})$.
The kernel of $P_j$ intersects $+1(s_{j-1})$ trivially by
Lemma~\ref{lem:triple}.
\end{proof}

\begin{corollary}
\label{cor:rank-structure}
In the staircase product $\Pi$, rank drops occur only at the $P_1$
steps. The total rank is $n!/2^d$ where $d$ is the number of
passes whose initial $P_1$ step causes a nontrivial rank reduction.
\end{corollary}


\section{First Block Surjectivity}

\begin{theorem}[First Block Surjectivity]
\label{thm:first-block}
After pass~1 (applying $P_1, P_2, \dots, P_{n-1}$ to $\CC[\Sym_n]$),
the image is exactly $+1(s_{n-1})$:
\[
  W_1 = P_{n-1} \cdots P_2 P_1(\CC[\Sym_n]) = +1(s_{n-1}).
\]
In particular, $\dim(W_1) = n!/2$.
\end{theorem}

\begin{proof}
After $P_1$: the image is $+1(s_1)$, which has dimension $n!/2$.
After $P_2$: the image lies in $+1(s_2)$. By Theorem~\ref{thm:injectivity},
$P_2$ is injective on $+1(s_1)$, so $\dim(P_2(+1(s_1))) = n!/2$.
Since $+1(s_2)$ also has dimension $n!/2$, we conclude
$P_2(+1(s_1)) = +1(s_2)$.

By induction: after $P_k$ (for $k = 2, \dots, n-1$), the image is all
of $+1(s_k)$, each of dimension $n!/2$. At the final step,
$W_1 = +1(s_{n-1})$.
\end{proof}


\section{Cascade Surjectivity}

\begin{lemma}[Cascade Surjectivity]
\label{lem:cascade}
If $|k - j| > 1$ (so $s_k$ and $s_j$ commute), then $P_k$ maps
$+1(s_{k-1}) \cap +1(s_j)$ bijectively onto $+1(s_k) \cap +1(s_j)$.
\end{lemma}

\begin{proof}
Since $s_k$ commutes with $s_j$, $P_k$ preserves $+1(s_j)$: if
$s_j v = v$, then $s_j(P_k v) = P_k(s_j v) = P_k v$.

By the Triple Lemma, $\ker(P_k) \cap +1(s_{k-1}) = \{0\}$, so $P_k$ is
injective on $+1(s_{k-1}) \cap +1(s_j)$. This space has dimension $n!/4$
(since $s_{k-1}$ and $s_j$ commute, both conditions independently halve).
The target $+1(s_k) \cap +1(s_j)$ also has dimension $n!/4$
(since $|k-j| > 1$ implies $s_k$ and $s_j$ commute). By injectivity plus
equal dimensions, the map is a bijection.
\end{proof}

\begin{remark}
More generally, if $S \subseteq \{1, \dots, n-1\}$ is a set of indices
with $|k-j| > 1$ for all $j \in S$ and $|k-j| > 1$, $|(k-1) - j| > 1$
for all $j \in S$, then $P_k$ maps
$+1(s_{k-1}) \cap \bigcap_{j \in S} +1(s_j)$ bijectively onto
$+1(s_k) \cap \bigcap_{j \in S} +1(s_j)$.
\end{remark}


\section{Trace Preservation}

\begin{lemma}[Trace Preservation]
\label{lem:trace}
Let $V \subseteq \CC[\Sym_n]$ be a subspace that is $s_j$-invariant
(i.e., $s_j(V) = V$). If $P_k$ commutes with $s_j$ (equivalently,
$|k-j| > 1$) and $P_k$ is injective on~$V$, then $P_k(V)$ is
$s_j$-invariant and
\[
  \tr(s_j|_{P_k(V)}) = \tr(s_j|_V).
\]
\end{lemma}

\begin{proof}
Since $P_k$ commutes with $s_j$, for $v \in V$:
$s_j(P_k v) = P_k(s_j v) \in P_k(V)$, so $P_k(V)$ is $s_j$-invariant.

Since $V$ is $s_j$-invariant, it decomposes as $V = V^+ \oplus V^-$
where $V^\pm = \pm 1(s_j) \cap V$. Then
$P_k(V) = P_k(V^+) \oplus P_k(V^-)$ (by injectivity of $P_k$ on~$V$
and the fact that $P_k$ preserves $\pm 1(s_j)$). Since $P_k$ is
injective, $\dim(P_k(V^\pm)) = \dim(V^\pm)$, so
\[
  \tr(s_j|_{P_k(V)}) = \dim(P_k(V^+)) - \dim(P_k(V^-))
  = \dim(V^+) - \dim(V^-) = \tr(s_j|_V). \qedhere
\]
\end{proof}


\section{The halving mechanism}

\subsection{Regularity of traces}

\begin{lemma}
\label{lem:reg-trace}
In the left regular representation of $\Sym_n$, for any non-identity
element $g \in \Sym_n$: $\tr(g) = 0$.
\end{lemma}

\begin{proof}
$g$ acts by $e_w \mapsto e_{gw}$. The trace is the number of fixed
points: $|\{w : gw = w\}| = 0$ since $g \ne e$.
\end{proof}

\begin{corollary}
\label{cor:trace-eigenspace}
For any set of pairwise commuting transpositions $\{s_{j_1}, \dots, s_{j_r}\}$
with $|j_a - j_b| > 1$ for all $a \ne b$, and any $s_i$ with
$|i - j_a| > 1$ for all $a$:
\[
  \tr\!\left(s_i \big|_{\bigcap_{a=1}^r +1(s_{j_a})}\right) = 0.
\]
\end{corollary}

\begin{proof}
The projector onto $\bigcap_a +1(s_{j_a})$ is $\prod_a P_{j_a}$
(commuting idempotents). The trace is
\[
  \tr\!\left(s_i \prod_a P_{j_a}\right)
  = 2^{-r} \sum_{T \subseteq \{1,\dots,r\}}
    \tr\!\left(s_i \prod_{a \in T} s_{j_a}\right).
\]
Each term $s_i \prod_{a \in T} s_{j_a}$ is a product of disjoint
transpositions including $s_i$, hence non-identity. By
Lemma~\ref{lem:reg-trace}, every term vanishes.
\end{proof}


\subsection{The first halving}

\begin{theorem}
\label{thm:halving-1}
For $n \ge 4$: $W_1 = +1(s_{n-1})$ is $s_1$-invariant, with
$\tr(s_1|_{W_1}) = 0$, so $P_1$ halves the image at the start
of pass~2:
\[
  \dim(P_1(W_1)) = \dim(W_1)/2 = n!/4.
\]
\end{theorem}

\begin{proof}
Since $|1 - (n-1)| = n-2 > 1$ for $n \ge 4$, $s_1$ commutes with
$s_{n-1}$, so $+1(s_{n-1})$ is $s_1$-invariant. By
Corollary~\ref{cor:trace-eigenspace} with $r = 1$ and $j_1 = n-1$,
$\tr(s_1|_{+1(s_{n-1})}) = 0$. Hence the $+1$ and $-1$ eigenspaces of
$s_1$ within $+1(s_{n-1})$ have equal dimension $n!/4$.
Applying $P_1$ kills the $-1$ eigenspace, giving
$\dim(P_1(W_1)) = n!/4$.
\end{proof}


\subsection{Image after pass~2}

\begin{theorem}
\label{thm:pass2}
For $n \ge 4$, the image after pass~2 is
\[
  W_2 = P_{n-2}\!\left(+1(s_{n-3}) \cap +1(s_{n-1})\right),
  \qquad \dim(W_2) = n!/4.
\]
Moreover, $W_2 \subseteq +1(s_{n-2})$.
\end{theorem}

\begin{proof}
Pass~2 applies $P_1, P_2, \dots, P_{n-2}$ to $W_1 = +1(s_{n-1})$.

After $P_1$: the image is $+1(s_1) \cap +1(s_{n-1})$, dimension $n!/4$
(by Theorem~\ref{thm:halving-1}).

For $k = 2, 3, \dots, n-3$: each $s_k$ commutes with $s_{n-1}$
(since $k \le n-3$ implies $|k - (n-1)| \ge 2$). By
Lemma~\ref{lem:cascade}, $P_k$ maps
$+1(s_{k-1}) \cap +1(s_{n-1})$ bijectively to
$+1(s_k) \cap +1(s_{n-1})$.

After $P_{n-3}$: the image is $+1(s_{n-3}) \cap +1(s_{n-1})$, dimension
$n!/4$.

Finally, $P_{n-2}$ is injective on this space (since
$+1(s_{n-3}) \cap +1(s_{n-1}) \subseteq +1(s_{n-3})$ and
$\ker(P_{n-2}) \cap +1(s_{n-3}) = \{0\}$ by the Triple Lemma).
So $W_2 = P_{n-2}(+1(s_{n-3}) \cap +1(s_{n-1}))$ has dimension $n!/4$,
and $W_2 \subseteq +1(s_{n-2})$.
\end{proof}


\subsection{The second halving}

\begin{theorem}
\label{thm:halving-2}
For $n \ge 6$: $W_2$ is $s_1$-invariant with $\tr(s_1|_{W_2}) = 0$,
so $P_1$ halves the image at the start of pass~3:
\[
  \dim(P_1(W_2)) = n!/8.
\]
\end{theorem}

\begin{proof}
The inner space $+1(s_{n-3}) \cap +1(s_{n-1})$ is $s_1$-invariant
(since $|1 - (n-3)| \ge 3$ and $|1 - (n-1)| \ge 5$ for $n \ge 6$),
and $\tr(s_1|_{+1(s_{n-3}) \cap +1(s_{n-1})}) = 0$ by
Corollary~\ref{cor:trace-eigenspace}.

The outer projection $P_{n-2}$ commutes with $s_1$ (since
$|1 - (n-2)| \ge 4$ for $n \ge 6$) and is injective on the inner
space (by the proof of Theorem~\ref{thm:pass2}). By the Trace
Preservation Lemma~\ref{lem:trace}:
\[
  \tr(s_1|_{W_2}) = \tr(s_1|_{+1(s_{n-3}) \cap +1(s_{n-1})}) = 0.
\]
So $W_2$ has $n!/8$ each of $+1$ and $-1$ eigenvectors for $s_1$,
and $P_1$ halves.
\end{proof}


\subsection{General halving pattern (verified)}

\begin{conjecture}[Halving Lemma]
\label{conj:halving}
For $j = 1, \dots, \lfloor n/2 \rfloor$: the $P_1$ step at the
start of pass~$j$ halves the image dimension. All subsequent $P_1$
steps (passes $\lfloor n/2 \rfloor + 1, \dots, n-1$) are injective.
\end{conjecture}

\noindent The above theorems prove this for the first two halvings
($n \ge 4$ and $n \ge 6$ respectively). The pattern continues:

\begin{remark}[Structure of later halvings]
Computational analysis (verified $n \le 8$) reveals:

\textbf{Image after pass~$j$:} $W_j \subseteq +1(s_{n-j})$,
with $\dim(W_j) = n!/2^{\min(j, \lfloor n/2 \rfloor)}$.

\textbf{$V_j$ formula:} Before halving $j{+}2$ (for
$j = 0, 1, \dots, \lfloor n/2 \rfloor - 2$), the image has the
form
\[
  W_{j+1} = P_{n-j-1}\!\left(
    +1(s_{n-j-2}) \cap +1(s_{n-j})
  \right) \cap (\text{$n!/2^{j+1}$-dim subspace}).
\]
More precisely, $W_{j+1}$ is obtained by applying $P_{n-j-1}$ to the
cascade output: the cascade $P_1, P_2, \dots, P_{n-j-2}$ maps
$+1(s_1) \cap W_j$ (after halving) through successive eigenspace
intersections, reaching $+1(s_{n-j-2}) \cap +1(s_{n-j})$ before the
final step $P_{n-j-1}$.

\textbf{$s_1$-invariance:} $W_{j+1}$ is $s_1$-invariant with
$\tr(s_1|_{W_{j+1}}) = 0$, for $j+1 \le \lfloor n/2 \rfloor - 1$.
This follows from:
\begin{enumerate}
  \item The ``inner'' intersection
    $+1(s_{n-j-2}) \cap +1(s_{n-j})$ is $s_1$-invariant with
    vanishing $s_1$-trace (by Corollary~\ref{cor:trace-eigenspace},
    provided $n - j - 2 \ge 3$, i.e., $j \le n - 5$).
  \item The ``outer'' projection $P_{n-j-1}$ commutes with $s_1$
    (provided $n - j - 1 \ge 3$, i.e., $j \le n - 4$).
  \item $P_{n-j-1}$ is injective on the inner space (by the
    Triple Lemma).
  \item Trace Preservation (Lemma~\ref{lem:trace}) propagates the
    vanishing trace through $P_{n-j-1}$.
\end{enumerate}
For $j+1 \le \lfloor n/2 \rfloor - 1$, the conditions
$j \le n - 5$ and $j \le n - 4$ are satisfied (since
$\lfloor n/2 \rfloor - 2 \le n - 5$ for $n \ge 6$).
\end{remark}

\noindent\textbf{Verification.} The halving pattern is confirmed for
$n = 2, 3, \dots, 8$:

\medskip
\begin{center}
\begin{tabular}{c|c|l}
$n$ & $\lfloor n/2 \rfloor$ & Rank after each $P_1$ step \\
\hline
2 & 1 & $2 \to 1$ \\
3 & 1 & $6 \to 3$ \\
4 & 2 & $24 \to 12 \to 6$ \\
5 & 2 & $120 \to 60 \to 30$ \\
6 & 3 & $720 \to 360 \to 180 \to 90$ \\
7 & 3 & $5040 \to 2520 \to 1260 \to 630$ \\
8 & 4 & $40320 \to 20160 \to 10080 \to 5040 \to 2520$
\end{tabular}
\end{center}

\medskip
\noindent Additionally, the cascade surjectivity formula
$W_2 = P_{n-2}(+1(s_{n-3}) \cap +1(s_{n-1}))$ is verified for
$n = 4, 5, 6$, with $s_1$-invariance and $\tr(s_1) = 0$ confirmed.


\section{The rank formula}

\begin{theorem}[Rank Formula]
\label{thm:rank}
Assuming Conjecture~\ref{conj:halving} (proved for the first two
halvings, verified for $n \le 8$):
\[
  \rk(\Pi) = \frac{n!}{2^{\lfloor n/2\rfloor}}.
\]
\end{theorem}

\begin{proof}
By Corollary~\ref{cor:rank-structure}, $\rk(\Pi) = n!/2^d$ where $d$ is
the number of halving $P_1$ steps. By Conjecture~\ref{conj:halving},
$d = \lfloor n/2 \rfloor$.
\end{proof}

\begin{remark}
The count $n!/2^{\lfloor n/2\rfloor}$ equals the number of
\emph{descent-paired permutations} of~$[n]$: those $w$ with
$w(2i{-}1) > w(2i)$ for all $i = 1, \dots, \lfloor n/2 \rfloor$
(OEIS A090932, Callan's characterization). Computational verification
($n \le 5$) confirms that $\{\Pi \cdot e_w : w \text{ descent-paired}\}$
is in fact a \emph{basis} for $\im(\Pi)$.
\end{remark}


\section{Eigenspace cascade tracking}

The cascade pattern through all steps reveals the following structure
(verified for $n \le 6$):

\begin{proposition}[Eigenspace tracking]
After step $P_k$ within pass~$j$ (i.e., after applying $P_1, \dots, P_k$
to $W_{j-1}$), the image is contained in $+1(s_k)$. Moreover, if
$k \le n - j - 1$, the image is also contained in $+1(s_{n-j+1})$
(the eigenspace condition ``carried'' from the previous pass).
At the last step $P_{n-j}$, the carried condition $+1(s_{n-j+1})$
is lost (since $s_{n-j}$ is adjacent to $s_{n-j+1}$), and the
image enters $+1(s_{n-j})$.
\end{proposition}

\noindent For $\Sym_6$, the full tracking is:
\begin{center}
\small
\begin{tabular}{r|l|l}
Step & Projection & Image fixed by \\
\hline
1 & $P_1$ & $s_1$ \\
2 & $P_2$ & $s_2$ \\
3 & $P_3$ & $s_3$ \\
4 & $P_4$ & $s_4$ \\
5 & $P_5$ & $s_5$ \quad (end of pass 1) \\
\hline
6 & $P_1$ & $s_1, s_5$ \quad (\textbf{halving}) \\
7 & $P_2$ & $s_2, s_5$ \\
8 & $P_3$ & $s_3, s_5$ \\
9 & $P_4$ & $s_4$ \quad ($s_5$ lost, end of pass 2) \\
\hline
10 & $P_1$ & $s_1, s_4$ \quad (\textbf{halving}) \\
11 & $P_2$ & $s_2, s_4$ \\
12 & $P_3$ & $s_3$ \quad ($s_4$ lost, end of pass 3) \\
\hline
13 & $P_1$ & $s_1$ \quad (\textbf{halving}) \\
14 & $P_2$ & $s_2$ \quad (end of pass 4) \\
\hline
15 & $P_1$ & $s_1$ \quad (no halving, end of pass 5) \\
\end{tabular}
\end{center}


\section{Remaining gaps}

\begin{enumerate}
  \item \textbf{General halving (passes $\ge 4$).} The proof of the
    first two halvings uses the explicit formula
    $W_2 = P_{n-2}(+1(s_{n-3}) \cap +1(s_{n-1}))$ to establish
    $s_1$-invariance and vanishing trace via Trace Preservation.
    Extending this to the third and higher halvings requires showing
    that the cascade $P_2, P_3, \dots$ (which includes $P_2$ that
    does \emph{not} commute with $s_1$) produces an $s_1$-invariant
    output. The difficulty: $P_2$ temporarily breaks $s_1$-invariance,
    and while $P_3, P_4, \dots$ (which commute with $s_1$) should
    ``restore'' it, a clean algebraic proof of this restoration
    is lacking.

  \item \textbf{Halving cessation.} After $\lfloor n/2 \rfloor$
    halvings, the image has no $-1(s_1)$ eigenvectors. This is
    verified computationally but not proved. The mechanism appears
    related to the ``carried'' eigenspace condition reaching
    indices too close to~1 for $s_1$-commutativity.

  \item \textbf{Image basis.} The conjecture that descent-paired
    permutations index an image basis is verified for $n \le 5$.
\end{enumerate}


\section{Proof summary and status}

\begin{center}
\begin{tabular}{l|l}
\textbf{Result} & \textbf{Status} \\
\hline
Triple Lemma (\ref{lem:triple}) & Proved \\
Within-Pass Injectivity (\ref{thm:injectivity}) & Proved \\
First Block Surjectivity (\ref{thm:first-block}) & Proved \\
Cascade Surjectivity (\ref{lem:cascade}) & Proved \\
Trace Preservation (\ref{lem:trace}) & Proved \\
Regularity of traces (\ref{lem:reg-trace}, \ref{cor:trace-eigenspace})
  & Proved \\
First halving (\ref{thm:halving-1}) & Proved ($n \ge 4$) \\
Image after pass~2 (\ref{thm:pass2}) & Proved ($n \ge 4$) \\
Second halving (\ref{thm:halving-2}) & Proved ($n \ge 6$) \\
General halving (\ref{conj:halving}) & Verified $n \le 8$ \\
Rank formula (\ref{thm:rank}) & Conditional on \ref{conj:halving} \\
Image basis & Verified $n \le 5$ \\
\end{tabular}
\end{center}


\begin{thebibliography}{9}
\bibitem{Galashin2020}
P.~Galashin, \emph{Symmetries of stochastic colored vertex models},
Ann.\ Probab.\ \textbf{49} (2021), no.~5, 2175--2219.

\bibitem{OEIS-A090932}
OEIS Foundation, \emph{A090932: $n!/2^{\lfloor n/2\rfloor}$},
\texttt{https://oeis.org/A090932}.
\end{thebibliography}

\end{document}
